\chapter*{Glosario de métricas y abreviaturas}
\addcontentsline{toc}{chapter}{Glosario de métricas y abreviaturas}

La memoria combina términos de aprendizaje automático, econometría y gestión de carteras. Este
glosario fija su significado antes de utilizarlos y evita redefinirlos en capítulos distintos.

\begin{description}[leftmargin=3.2cm,style=nextline,itemsep=0.45em]
  \item[Alfa] Retorno de la cartera que excede al benchmark. En esta memoria puede designar el
  exceso realizado o, cuando se indica expresamente, el exceso esperado que gobierna una orden.
  \item[Amplitud] Número aproximado de decisiones independientes sobre las que puede expresarse una
  señal. Doce cohortes mensuales solapadas a horizonte anual no equivalen a doce observaciones
  independientes.
  \item[Benchmark] Referencia contra la que se mide la cartera. SPY se usa como representación
  invertible del S\&P 500, pero nunca puede ser una posición de la estrategia.
  \item[Bootstrap por bloques] Remuestreo que conserva tramos consecutivos de cohortes en lugar de
  extraerlas una a una. Con etiquetas solapadas a doce meses, romper esa dependencia produciría
  intervalos artificialmente estrechos.
  \item[Coeficiente de transferencia] Fracción de la señal disponible que llega efectivamente a la
  cartera. Es el término de la ley fundamental de la gestión activa que las restricciones
  \textit{long-only}, la concentración y los costes reducen.
  \item[CAGR] Tasa de crecimiento anual compuesta. Resume el crecimiento geométrico del patrimonio
  y no informa por sí sola de su volatilidad ni de su recorrido.
  \item[Cohorte] Sección transversal de acciones puntuada en una misma fecha y evaluada contra el
  retorno que se realiza al final del horizonte declarado. Está \emph{cerrada} cuando ese horizonte
  ya ha transcurrido y su retorno es observable; solo entonces puede entrar en un entrenamiento o en
  un diagnóstico.
  \item[DSR] \textit{Deflated Sharpe Ratio}. Probabilidad de que el Sharpe observado supere lo
  esperable por azar después de considerar multiplicidad, asimetría y curtosis.
  \item[Drawdown] Caída desde un máximo previo del patrimonio. El máximo drawdown es la peor de esas
  pérdidas durante una ventana.
  \item[Empate técnico] Situación en la que dos alternativas difieren menos de 0,002 en ventaja
  pareada. La regla declarada de antemano la resuelve por simplicidad, no por magnitud, de modo que
  su resultado no debe presentarse como un hallazgo empírico.
  \item[HHI] Índice de Herfindahl--Hirschman aplicado a los pesos de cartera. Aumenta cuando el
  patrimonio se concentra en menos posiciones.
  \item[IC-IR] Media temporal del Rank-IC dividida por su desviación estándar. Mide estabilidad de
  la capacidad de ordenar, no rentabilidad de cartera.
  \item[Information Ratio (IR)] Media del exceso de retorno dividida por su volatilidad. A
  diferencia del IC-IR, pertenece a la capa económica y depende de la cartera.
  \item[No inferioridad] Puerta alternativa a la dominancia: un candidato es admisible si el límite
  inferior de su intervalo \textit{bootstrap} pareado al 90\,\% supera el margen declarado. Prueba
  que no es peor, no que sea mejor.
  \item[OOS] \textit{Out of sample}. Observaciones posteriores a las usadas para ajustar el modelo.
  En esta memoria también se distingue dentro de OOS entre selección y reserva.
  \item[Placebo de etiqueta] Ejecución completa del sistema con el retorno futuro barajado dentro de
  cada cohorte. Si la maquinaria fabricase señal por sí misma, el placebo la mostraría.
  \item[P95] Percentil 95. En capacidad es la participación sobre volumen que no supera el 95\,\%
  de las órdenes; evita que un único extremo gobierne la estimación.
  \item[PIT] \textit{Point in time}. Información asociada a la fecha en que era observable, no a
  la fecha económica a la que se refiere ni a una versión revisada conocida después.
  \item[Rank-IC] Correlación de Spearman entre el ranking emitido y el retorno futuro excedente
  dentro de una cohorte. Es la métrica que selecciona el modelo.
  \item[Reserva] Periodo 2025--2026 aislado de toda selección. Sus pocas cohortes cerradas permiten
  una comprobación preliminar, no una estimación precisa de generalización.
  \item[Slippage] Diferencia adversa entre el precio de referencia y el de ejecución. Se suma a la
  comisión para calcular el coste de una operación.
  \item[Snapshot] Fecha en la que el sistema reconstruye el universo observable, emite rankings y
  decide la cartera.
  \item[Suma cero] Propiedad aritmética del problema de batir a un índice: antes de costes, el
  conjunto de la gestión activa obtiene exactamente el rendimiento del mercado, de modo que el
  exceso de un participante es el defecto de otro.
  \item[Turnover] Suma anualizada de cambios absolutos de peso. Una cifra de 1,35 equivale a mover
  nocionalmente un 135\,\% del patrimonio en un año, no a cerrar la cartera 1,35 veces.
  \item[Ventila] Cada uno de los veinte tramos en que se divide una cohorte ordenada por
  \texttt{meta\_rank}. Son la unidad sobre la que se estima la curva que traduce rango en alfa
  esperada.
  \item[Ventana de selección] Cohortes de 2015--2024 sobre las que se comparan configuraciones. La
  era reservada queda fuera de todas esas decisiones.
  \item[Walk-forward] Esquema que avanza cronológicamente: ajusta con historia disponible, predice
  la fecha siguiente y repite sin permitir que el futuro vuelva al entrenamiento.
\end{description}

